\documentclass[11pt]{article}

\usepackage[T1]{fontenc}
\usepackage[utf8]{inputenc}
\usepackage{microtype}
\usepackage[a4paper,margin=30mm]{geometry}
\usepackage{amsmath,amssymb,amsthm,mathtools}
\usepackage{lmodern}
\usepackage[hidelinks]{hyperref}
\hypersetup{
  pdftitle={The Subfield-Coefficient Root Distribution for a Seventh-Degree Unit-Circle Polynomial},
  pdfauthor={Xiaojian Zeng},
  pdfsubject={A root-count distribution over finite fields of characteristic two},
  pdfkeywords={finite fields, unit circle, additive characters, Artin--Schreier curve}
}

\newtheorem{theorem}{Theorem}
\newtheorem{lemma}[theorem]{Lemma}
\newtheorem{remark}[theorem]{Remark}

\newcommand{\F}{\mathbb F}
\newcommand{\Tr}{\operatorname{Tr}}

\title{The Subfield-Coefficient Root Distribution for a\\
Seventh-Degree Unit-Circle Polynomial}
\author{Xiaojian Zeng\\
\small PBC School of Finance, Tsinghua University, Beijing 100190, China\\
\small \texttt{xiaojianzeng@foxmail.com}}
\date{September 9, 2026}

\begin{document}
\maketitle

\begin{abstract}
Let $q=2^m$ and let $U_{q+1}$ be the norm-one subgroup of
$\F_{q^2}^{\times}$.  For $a\in\F_q$, write $N_m(a)$ for the number of
distinct roots in $U_{q+1}$ of
\[
 x^7+a x^4+a x^3+1.
\]
We determine the distribution of $N_m(a)$ for every $m\geq1$.  If
$c_1=1$, $c_2=-3$, $c_m=c_{m-1}-2c_{m-2}$, and
$d_m=-2+(-1)^m+c_m$, then the respective frequencies of five, three,
and one roots are
\[
 \frac{2^m+d_m}{8},\qquad
 \frac{2^m-d_m}{4},\qquad
 \frac{5\cdot2^m+d_m}{8}.
\]
The proof reduces the distribution to an ordered collision count, cancels
two additive-character sums, and identifies the remaining sum with the
point count of the genus-one curve
$(t^2+t)(w^2+w+1)=1$ over $\F_2$.  All exceptional fibres are included
explicitly.
\end{abstract}

\section{Statement and context}

Put $q=2^m$ and
\[
 U_{q+1}=\{x\in\F_{q^2}^{\times}:x^{q+1}=1\}.
\]
For $a\in\F_q$, let
\[
 N_m(a)=\#\{x\in U_{q+1}:x^7+a x^4+a x^3+1=0\},
 \qquad
 M_j(m)=\#\{a\in\F_q:N_m(a)=j\}.
\]
Roots are counted without multiplicity.  The following is the result proved
in this note.

\begin{theorem}\label{thm:main}
For every integer $m\geq1$, define
\[
 c_1=1,\qquad c_2=-3,\qquad
 c_m=c_{m-1}-2c_{m-2}\quad(m\geq3),
\]
and $d_m=-2+(-1)^m+c_m$.  Then
\[
 M_5(m)=\frac{2^m+d_m}{8},\qquad
 M_3(m)=\frac{2^m-d_m}{4},\qquad
 M_1(m)=\frac{5\cdot2^m+d_m}{8}.
\]
\end{theorem}

Xiong and Yan study the ambient family
$x^7+a x^4+\bar a x^3+1$ with $a\in\F_{q^2}$ and determine its complete
root-count distribution on $U_{q+1}$ \cite{XiongYan2026}.  The specialization
considered here imposes $a=\bar a$, hence $a\in\F_q$.  An ambient
$q^2$-parameter distribution does not determine this conditional
$q$-parameter distribution; Theorem~\ref{thm:main} evaluates the latter.
No priority claim is made.

Every algebraic reduction below is specific to characteristic two.  The
importance of recording the ground-field hypothesis is also visible in other
algebraic problems: for example, a canonical Bézout map studied by Zeng is
generically reversible in characteristic zero but loses information on an
explicit positive-characteristic family \cite{Zeng2026Bezout}.  That paper is
only a comparison concerning characteristic dependence and supplies no input
to the present proof.

We use the absolute trace $\Tr=\Tr_{\F_q/\F_2}$ and its canonical additive
character $\psi(u)=(-1)^{\Tr(u)}$.  Standard background on trace criteria and
additive characters over finite fields may be found in
\cite{LidlNiederreiter1996}.

\section{Reduction to collisions}

Define
\[
 S=\{y\in\F_q^{\times}:\Tr(1/y)=1\}.
\]

\begin{lemma}\label{lem:unit-circle}
The map $x\mapsto x+x^{-1}$ is two-to-one from
$U_{q+1}\setminus\{1\}$ onto $S$.  Moreover, if
\[
 R(a)=\#\{y\in S:y^3+y^2+1=a\},
\]
then $N_m(a)=1+2R(a)$ and $R(a)\leq2$.
\end{lemma}

\begin{proof}
In characteristic two,
\begin{equation}\label{eq:factor}
 x^7+a x^4+a x^3+1
 =(x+1)\bigl(x^6+x^5+x^4+(a+1)x^3+x^2+x+1\bigr).
\end{equation}
Thus $x=1$ is always a root.  If
$x\in U_{q+1}\setminus\{1\}$ and $y=x+x^{-1}$, then
$y\in\F_q^{\times}$.  Dividing the degree-six factor in
\eqref{eq:factor} by $x^3$ gives
\begin{equation}\label{eq:cubic}
 y^3+y^2=a+1.
\end{equation}

Conversely, the inverse fibre over $y\in\F_q^{\times}$ consists of the
roots of $X^2+yX+1$.  After scaling, the usual characteristic-two quadratic
criterion says that this polynomial is irreducible over $\F_q$ precisely when
\[
 \Tr(1/y^2)=\Tr(1/y)=1.
\]
In that event its two roots are Frobenius conjugates with product one, so
they lie in $U_{q+1}$.  This proves the asserted two-to-one map and, together
with \eqref{eq:cubic}, the identity $N_m(a)=1+2R(a)$.

It remains to bound the fibre size.  If three distinct members
$y_1,y_2,y_3$ of $S$ were roots of
$Y^3+Y^2+(a+1)$, their product would be nonzero and Vieta's relations would
give
\[
 \frac1{y_1}+\frac1{y_2}+\frac1{y_3}=0.
\]
Taking absolute traces would give $0=1+1+1=1$ in $\F_2$, a contradiction.
When $a+1=0$, the polynomial is $Y^2(Y+1)$ and has only one distinct
nonzero root.  Hence $R(a)\leq2$ in all cases.
\end{proof}

Let $D_m$ be the number of ordered pairs $(y,z)\in S^2$ such that
$y\ne z$ and $y^3+y^2=z^3+z^2$.  Since every fibre of $R$ has size at
most two,
\begin{equation}\label{eq:D-fibres}
 D_m=2\#\{a\in\F_q:R(a)=2\}.
\end{equation}
Cancellation of $y+z\ne0$ shows that the collision equation is
\begin{equation}\label{eq:collision}
 y^2+yz+z^2+y+z=0.
\end{equation}

\section{The character sum}

Set $T=\F_q\setminus\{0,1\}$.  In \eqref{eq:collision}, put $t=z/y$.
For an actual collision,
\begin{equation}\label{eq:param}
 y=\frac{1+t}{1+t+t^2},\qquad
 z=\frac{t(1+t)}{1+t+t^2},
\end{equation}
where $t\in T$ and $1+t+t^2\ne0$.  The two reciprocal trace arguments are
\begin{equation}\label{eq:AB}
 A(t)=\frac1y=t+\frac1{t+1},\qquad
 B(t)=\frac1z=1+\frac1t+\frac1{t+1}.
\end{equation}
If $1+t+t^2=0$, both expressions in \eqref{eq:AB} are zero.  Such a
parameter contributes zero to the simultaneous trace-one indicator, so it
may be included without changing the count.  Therefore
\begin{equation}\label{eq:indicator}
 D_m=\frac14\sum_{t\in T}
       (1-\psi(A(t)))(1-\psi(B(t))).
\end{equation}

Translation $t\mapsto t+1$ permutes $T$, and
\[
 A(t+1)=t+1+\frac1t=A(t)+B(t).
\]
Consequently, in the expansion of \eqref{eq:indicator}, the $A$-sum
cancels the $(A+B)$-sum.  Writing
\begin{equation}\label{eq:sigma}
 \Sigma_m=\sum_{t\in T}
 \psi\!\left(1+\frac1{t(t+1)}\right),
\end{equation}
we obtain the exact identity
\begin{equation}\label{eq:D-sigma}
 D_m=\frac{q-2-\Sigma_m}{4}.
\end{equation}

\section{A genus-one point count}

Consider the projective closure $E$ in
$\mathbb P^1_t\times\mathbb P^1_w$ of
\begin{equation}\label{eq:curve}
 (t^2+t)(w^2+w+1)=1
\end{equation}
over $\F_2$.  We verify the geometry needed for its point count rather than
leaving exceptional fibres implicit.

\begin{lemma}\label{lem:curve}
The curve $E$ is smooth and geometrically integral of genus one.  It has
exactly two $\F_2$-rational points, and
\begin{equation}\label{eq:E-count}
 \#E(\F_{2^m})=2^m+1-c_m,
\end{equation}
where $(c_m)$ is the recurrence in Theorem~\ref{thm:main}.
\end{lemma}

\begin{proof}
On the affine chart, put $P=t^2+t$ and $Q=w^2+w+1$.  The two partial
derivatives of $PQ+1$ are $Q$ and $P$, and they cannot vanish on the curve.
Near $t=\infty$, with $s=1/t$, the equation becomes
\[
 (1+s)Q=s^2;
\]
at $s=0$ its derivative with respect to $w$ is one.  Near $w=\infty$,
with $v=1/w$, it becomes
\[
 P(1+v+v^2)=v^2;
\]
at $v=0$ its derivative with respect to $t$ is one.  The chart where both
coordinates are infinite contains no point of $E$.  Thus $E$ is smooth.

The Artin--Schreier presentation
\begin{equation}\label{eq:AS}
 w^2+w=1+\frac1{t(t+1)}
\end{equation}
also proves geometric integrality.  Indeed, its right-hand side has simple
poles at $t=0,1$, whereas every pole of $h^2+h$ has even order.  Hence the
right-hand side is not of that form over $\overline{\F}_2(t)$, so the
Artin--Schreier polynomial is irreducible.  Since $E$ is a smooth integral
curve of bidegree $(2,2)$, the adjunction formula gives
$g(E)=(2-1)(2-1)=1$.  These Artin--Schreier and function-field facts are
standard; see \cite{Stichtenoth2009}.

The two points of $E(\F_2)$ are the points with $w=\infty$ over $t=0$ and
$t=1$.  There is no point above $t=\infty$, because that fibre requires
$w^2+w+1=0$, which has no root in $\F_2$.  Hence
$\#E(\F_2)=2$, and the Frobenius trace is $2+1-2=1$.  The Frobenius
polynomial is therefore
\[
 X^2-X+2.
\]
If its roots are $\alpha,\beta$, then the point-count formula for curves over
finite fields \cite{Weil1948} gives
\[
 \#E(\F_{2^m})=2^m+1-(\alpha^m+\beta^m).
\]
Now $\alpha+\beta=1$ and $\alpha\beta=2$, so
$\alpha^m+\beta^m$ has initial values $1,-3$ and satisfies
$c_m=c_{m-1}-2c_{m-2}$.  This proves \eqref{eq:E-count}.
\end{proof}

We count the same points directly from \eqref{eq:AS}.  For each finite
$t\in T$, the number of $w$ is
\[
 1+\psi\!\left(1+\frac1{t(t+1)}\right).
\]
The poles $t=0,1$ contribute the two rational points with $w=\infty$.
Above $t=\infty$, the equation is $w^2+w+1=0$; this contributes two
points if $m$ is even and none if $m$ is odd, namely $1+(-1)^m$ points.
Thus
\begin{align}
 \#E(\F_q)
  &=(q-2)+\Sigma_m+2+1+(-1)^m \notag\\
  &=q+1+(-1)^m+\Sigma_m.\label{eq:fibres}
\end{align}
Comparing \eqref{eq:E-count} and \eqref{eq:fibres} yields
\begin{equation}\label{eq:sigma-value}
 \Sigma_m=-(-1)^m-c_m.
\end{equation}
Together with \eqref{eq:D-sigma}, this gives
\begin{equation}\label{eq:D-value}
 D_m=\frac{q-2+(-1)^m+c_m}{4}=\frac{q+d_m}{4}.
\end{equation}

\section{Recovery of the frequencies}

Let $n_r=\#\{a\in\F_q:R(a)=r\}$ for $r=0,1,2$.  Inversion is a
bijection of $\F_q^{\times}$, and exactly $q/2$ field elements have absolute
trace one.  Consequently,
\[
 n_0+n_1+n_2=q,\qquad n_1+2n_2=q/2.
\]
Equations \eqref{eq:D-fibres} and \eqref{eq:D-value} now give
\[
 n_2=\frac{q+d_m}{8},\qquad
 n_1=\frac{q-d_m}{4},\qquad
 n_0=\frac{5q+d_m}{8}.
\]
Finally $N_m(a)=1+2R(a)$, so $(n_2,n_1,n_0)$ is
$(M_5(m),M_3(m),M_1(m))$.  This proves Theorem~\ref{thm:main}, including
$m=1$.

\section*{Computational and AI assistance}

OpenAI Codex assisted with proof development, exact finite-field checks,
literature searches, and preparation of the manuscript.  Package-free exact
verifiers independently check the reduced fibre distribution, the collision
count, the character sum, and the curve point count for every
$1\leq m\leq15$.  These finite calculations are corroborating diagnostics;
the all-$m$ conclusion rests on the symbolic argument above.  The verifier
source and captured output are available with the accompanying project
materials.  Responsibility for any use of this manuscript remains with the
named author.

\bibliographystyle{plain}
\bibliography{references}

\end{document}
